%%%%%%%%%%%%%%%%%%%%%%%%%%%%%%%%%%%%%%%%%%%%%%%%%%%%%%%%%%%%%%%%%%%%%%%%%%%
%% cell_reweighting_report.tex
%%
%% Physics report: Photon shower shape corrections using cell-level energy
%% reweighting with Run 3 data.
%%
%% This documents the CURRENT pipeline (M1 flat shift + M2 shift+stretch).
%% The older investigation of the DeltaR-matched TProfile method is in
%% data_mc_reweighting_report.tex (kept for reference).
%%
%% Compile:
%%   cd report/
%%   pdflatex cell_reweighting_report.tex   (x2 for references)
%%
%% Prerequisites:
%%   - Run the pipeline:  cd scripts/data_mc && bash run.sh --batch
%%   - Run chi2 extraction:  root -l -b -q 'extract_chi2.C'
%%   This generates chi2_yields.tex, chi2_summary.tex, chi2_tables.tex.
%%%%%%%%%%%%%%%%%%%%%%%%%%%%%%%%%%%%%%%%%%%%%%%%%%%%%%%%%%%%%%%%%%%%%%%%%%%

\documentclass[11pt,a4paper]{article}

\usepackage[utf8]{inputenc}
\usepackage[T1]{fontenc}
\usepackage[margin=2.5cm]{geometry}
\usepackage{graphicx}
\usepackage{amsmath,amssymb}
\usepackage{booktabs,longtable}
\usepackage{caption,subcaption}
\usepackage{float,xcolor}
\usepackage{hyperref}
\usepackage{siunitx}
\usepackage{array}

\hypersetup{
  colorlinks=true,
  linkcolor=blue!50!black,
  citecolor=blue!50!black,
  urlcolor=blue!50!black
}

% Graphics path: output directory (relative to report/)
\graphicspath{{../output/cell_energy_reweighting_Francisco_method/data24/}}

% Custom commands
\newcommand{\Reta}{\ensuremath{R_{\eta}}}
\newcommand{\Rphi}{\ensuremath{R_{\phi}}}
\newcommand{\weta}{\ensuremath{w_{\eta 2}}}
\newcommand{\abseta}{\ensuremath{|\eta|}}
\newcommand{\chisq}{\ensuremath{\chi^2 / n_{\mathrm{df}}}}
\newcommand{\pt}{\ensuremath{p_{\mathrm{T}}}}
\newcommand{\Etot}{\ensuremath{E_{\mathrm{total}}}}

\title{Photon Shower Shape Corrections Using Cell-Level Energy Reweighting\\
       with ATLAS Run\,3 Data}
\author{Marta Fernandes da Silva}
\date{April 2026}

\begin{document}
\maketitle

\begin{abstract}
Electromagnetic shower shapes (\Reta, \Rphi, \weta) are critical inputs to
photon identification in ATLAS.  Residual data--MC mismodelling of the
calorimeter cell energy deposits degrades the accuracy of these
discriminating variables.  This note presents a cell-level energy
reweighting approach applied to Layer\,2 of the electromagnetic
calorimeter using $Z \to \ell\ell\gamma$ final-state-radiation events from
ATLAS Run\,3 data (2024) and MC23e simulation.

Two correction methods are studied.  Method~1 (M1, flat shift) adjusts
each cell energy by the mean data--MC fraction difference.  Method~2
(M2, shift+stretch) additionally matches the variance by applying a
linear transformation derived from Francisco Almeida's
\texttt{photoncellbasedrw} framework.  Both methods are evaluated across
14~$\abseta$ bins, three decay channels ($Z \to ee\gamma$,
$Z \to \mu\mu\gamma$, combined), and three photon conversion
scenarios (unconverted, converted, inclusive).

M1 provides robust corrections for \Reta\ and \Rphi\ with minimal
residual $\chisq$ and exact energy conservation by construction.  M2
further improves \weta\ by matching the cell-fraction variance.
Directions for future work include differential \pt-dependent
corrections and machine-learning approaches such as normalising flows.
\end{abstract}

\tableofcontents
\newpage


%% ====================================================================
\section{Introduction}
\label{sec:intro}
%% ====================================================================

Photon identification in ATLAS relies on shower shape variables computed
from the energy deposits in the electromagnetic calorimeter.  In
particular, \Reta, \Rphi, and \weta\ are built from the Layer\,2
cell energies in a $7 \times 11$ ($\eta \times \phi$) cluster and are
among the most discriminating variables for separating prompt photons
from hadronic fakes.

Residual mismodelling of these variables in Monte Carlo simulation is
addressed by the ATLAS ``fudge factor'' framework, which applies
empirical corrections to the branch-level shower shape values.  An
alternative approach, pioneered by Francisco Almeida in Run\,2
\cite{francisco}, operates at the cell-energy level: the individual cell
energies are corrected before recomputing the shower shapes, thereby
preserving the physical correlations between cells.

This note documents the Run\,3 implementation of the cell-level
reweighting approach.  Two methods are studied:
\begin{itemize}
  \item \textbf{Method~1 (M1)}: a flat shift per cell, matching the
        population mean of the cell energy fraction.
  \item \textbf{Method~2 (M2)}: a shift+stretch transformation, matching
        both the mean and the variance, following Francisco Almeida's
        \texttt{photoncellbasedrw} method \cite{francisco}.
\end{itemize}
The analysis uses $Z \to \ell\ell\gamma$ events selected from ATLAS 2024
data and MC23e simulation, processed through the NTupleMaker framework.


%% ====================================================================
\section{Data and MC Samples}
\label{sec:samples}
%% ====================================================================

The analysis uses ATLAS data collected in 2024 and the corresponding
MC23e simulation campaign.  Two channels are considered:
\begin{itemize}
  \item $Z \to ee\gamma$: data from DAOD\_EGAM3 (\texttt{egam3.root}),
        MC sample DSID~700770.
  \item $Z \to \mu\mu\gamma$: data from DAOD\_EGAM4 (\texttt{egam4.root}),
        MC sample DSID~700771.
\end{itemize}
The combined channel ($Z \to \ell\ell\gamma$) merges both data and MC
samples.  Ntuples are produced by the NTupleMaker package running in
Athena~25.0.40.

MC events are weighted by $w = w_{\text{MC}} \times w_\mu \times \sigma$,
where $w_{\text{MC}}$ is the generator weight, $w_\mu$ the pile-up
reweighting factor, and $\sigma$ the cross-section weight.

% Auto-generated by extract_chi2.C — do not edit by hand
\begin{table}[htbp]
\centering
\caption{Event yields after all selection cuts for each channel and conversion scenario.  Data counts are unweighted; MC counts are unweighted (the MC event weight $w = w_{\text{MC}} \times w_\mu \times \sigma$ is applied to histograms but not to the event count).}
\label{tab:yields}
\begin{tabular}{llrr}
\toprule
Channel & Scenario & Data events & MC events \\
\midrule
$Z \to \ell\ell\gamma$ & Unconverted & 205309 & 2793063 \\
$Z \to \ell\ell\gamma$ & Converted & 45532 & 756176 \\
$Z \to \ell\ell\gamma$ & Inclusive & 250841 & 3549239 \\
\bottomrule
\end{tabular}
\end{table}


%% ====================================================================
\section{Event Selection}
\label{sec:selection}
%% ====================================================================

Events are selected in two stages: a preselection applied at the
NTupleMaker level during ntuple production, and an offline selection
applied by the analysis pipeline.

\subsection{NTupleMaker-level preselection}
\label{sec:preselection}

The NTupleMaker applies Good Runs List (GRL) filtering, trigger
matching, and basic lepton quality requirements.  Photon candidates are
required to have a complete $7 \times 11$ Layer\,2 cell cluster
(77~cells) stored in the ntuple output.

\subsection{Offline selection}
\label{sec:offline-selection}

The analysis pipeline applies the following offline cuts:

\begin{table}[htbp]
\centering
\caption{Summary of offline selection requirements.}
\label{tab:selection}
\begin{tabular}{ll}
\toprule
Requirement & Value \\
\midrule
Photon \pt         & $> 10$~GeV \\
Photon $\abseta$   & $< 2.37$, excluding $1.37 \leq \abseta < 1.52$ \\
Dilepton mass      & $40 < m_{\ell\ell} < 83$~GeV \\
Three-body mass    & $80 < m_{\ell\ell\gamma} < 100$~GeV \\
$\Delta R(\ell, \gamma)$ & $> 0.4$ (both leptons) \\
Cluster size       & exactly 77 cells \\
Central cell       & must carry the highest energy (safety check$^{*}$) \\
Total cluster energy & $> 0$ \\
\bottomrule
\end{tabular}
\smallskip

\noindent$^{*}$The NTupleMaker constructs the $7\times11$ cluster
centred on the hottest Layer\,2 cell, so this requirement is satisfied
by construction; it is imposed as a redundant safety check.
MC truth matching is available in the ntuple (\texttt{requireTruthMC}
flag in \texttt{config.h}) but is \textbf{not applied} in the current
configuration, consistent with Francisco's reference implementation.
\end{table}

The dilepton and three-body mass windows select the final-state-radiation
(FSR) topology $Z \to \ell\ell\gamma$.  The central-cell requirement
is a safety net, since the NTupleMaker already centres the cluster on
the hottest cell.

\subsection{Conversion scenarios}
\label{sec:conversion-scenarios}

Three conversion scenarios are studied:
\begin{itemize}
  \item \textbf{Baseline} (unconverted): only unconverted photons
        ($\texttt{isConv} = \texttt{false}$), with loose isolation.
  \item \textbf{Converted}: only converted photons
        ($\texttt{isConv} = \texttt{true}$), with loose isolation.
  \item \textbf{All conversions}: no conversion requirement, with loose
        isolation.
\end{itemize}


%% ====================================================================
\section{Shower Shape Variables}
\label{sec:shower-shapes}
%% ====================================================================

The three shower shape variables are computed from the $7 \times 11$
Layer\,2 cell energy cluster centred on the cell with the highest energy
(index~38 in the row-major ordering).  The cell granularity is
$\Delta\eta = 0.025$ and $\Delta\phi \approx 0.0245$.

\begin{description}
  \item[\Reta] Ratio of the energy in a $3 \times 7$ ($\eta \times \phi$)
    window to the energy in the full $7 \times 7$ window:
    \begin{equation}
      \Reta = \frac{E_{3 \times 7}}{E_{7 \times 7}} \,.
    \end{equation}

  \item[\Rphi] Ratio of the energy in a $3 \times 3$ window to the energy
    in a $3 \times 7$ window:
    \begin{equation}
      \Rphi = \frac{E_{3 \times 3}}{E_{3 \times 7}} \,.
    \end{equation}

  \item[\weta] Energy-weighted $\eta$ width in a $3 \times 5$ window
    centred on the hottest cell:
    \begin{equation}
      \weta = \sqrt{\frac{\sum_i E_i \, (\eta_i - \bar\eta)^2}{\sum_i E_i}}
      \,,\qquad
      \bar\eta = \frac{\sum_i E_i \, \eta_i}{\sum_i E_i} \,,
    \end{equation}
    where the sum runs over the $3 \times 5$ cells.  A position-dependent
    correction is applied following the ATLAS convention.
\end{description}

The energy windows $E_{n \times m}$ are computed from the cell energies
using the cell $\eta$ and $\phi$ positions to identify which cells fall
within the window, following the implementation in \texttt{config.h}.

\paragraph{Validation.}
The cell-computed shower shapes are validated against the values stored
in the ntuple branches (which are the official ATLAS reconstruction
values).  Figure~\ref{fig:computed-vs-stored} shows excellent agreement,
confirming that the cell-level computation faithfully reproduces the
standard shower shapes.

\begin{figure}[htbp]
\centering
\includegraphics[width=0.92\textwidth]{llgamma/baseline/plots/computed_vs_stored.pdf}
\caption{Validation: cell-computed shower shapes versus stored ntuple
  branch values for the combined $Z \to \ell\ell\gamma$ channel
  (unconverted photons).}
\label{fig:computed-vs-stored}
\end{figure}


%% ====================================================================
\section{Correction Methods}
\label{sec:methods}
%% ====================================================================

Both correction methods operate on the cell energy fractions
$f_k = E_k / \Etot$, where $E_k$ is the energy of cell~$k$ and
$\Etot = \sum_{k=1}^{77} E_k$ is the total cluster energy.
Corrections are derived independently for each of 14~$\abseta$ bins
(Section~\ref{sec:selection}).

\subsection{Method~1: flat shift}
\label{sec:m1}

For each cell~$k$ and $\abseta$ bin, the correction is defined as the
difference in population-mean fractions between data and MC:
\begin{equation}
  \Delta_k = \langle f_{\text{data}} \rangle_k
           - \langle f_{\text{MC}} \rangle_k \,.
  \label{eq:m1-delta}
\end{equation}
The corrected cell energy is then
\begin{equation}
  E'_k = E_k + \Delta_k \times \Etot \,.
  \label{eq:m1-apply}
\end{equation}

\paragraph{Energy conservation.}
Since $\sum_k \langle f_{\text{data}} \rangle_k = \sum_k \langle
f_{\text{MC}} \rangle_k = 1$, the sum of corrections vanishes:
$\sum_k \Delta_k = 0$.  Therefore $\sum_k E'_k = \Etot$ exactly, with
no need for renormalisation.

M1 is simple and robust: it corrects the mean cell profile but does not
modify the variance of the cell fraction distribution.

\subsection{Method~2: shift+stretch (Francisco's method)}
\label{sec:m2}

Method~2 additionally matches the standard deviation of the cell
fraction distribution.  For each cell~$k$, define:
\begin{align}
  \text{stretch}_k &= \frac{\sigma_{\text{data},k}}{\sigma_{\text{MC},k}}
  \,,\label{eq:m2-stretch}\\
  \text{shift}_k &= \mu_{\text{data},k}
                   - \text{stretch}_k \times \mu_{\text{MC},k}
  \,,\label{eq:m2-shift}
\end{align}
where $\mu$ and $\sigma$ denote the population mean and standard
deviation of the cell fraction $f_k$ in data and MC respectively.
The corrected cell energy is
\begin{equation}
  E'_k = \Etot \times \text{shift}_k
       + \text{stretch}_k \times E_k \,.
  \label{eq:m2-apply}
\end{equation}

This is a linear transformation of the cell energy.  At the population
level:
\begin{itemize}
  \item The mean is matched:
    $\langle E'_k / \Etot \rangle = \text{stretch}_k \cdot \mu_{\text{MC},k}
    + \text{shift}_k = \mu_{\text{data},k}$.
  \item The standard deviation is matched:
    $\sigma'_k = \text{stretch}_k \cdot \sigma_{\text{MC},k}
    = \sigma_{\text{data},k}$.
\end{itemize}

\paragraph{Energy conservation.}
Unlike M1, M2 does not automatically conserve total energy.  A
rescaling step is applied after correction:
\begin{equation}
  E'_k \leftarrow E'_k \times \frac{\Etot}{\sum_{j} E'_j} \,.
  \label{eq:m2-rescale}
\end{equation}

This method corresponds to Francisco Almeida's
\texttt{photoncellbasedrw} approach \cite{francisco}, adapted for
Run\,3 data.


%% ====================================================================
\section{Results: Baseline (Unconverted Photons)}
\label{sec:results-baseline}
%% ====================================================================

This section presents the main results for the combined channel
($Z \to \ell\ell\gamma$) with unconverted photons (baseline scenario).

\subsection{Shower shape comparisons}
\label{sec:results-shower}

Figures~\ref{fig:rew-reta}--\ref{fig:rew-weta2} show the cell-computed
shower shapes for data, uncorrected MC, and MC after M1 and M2
corrections.  Each PDF contains one page per $\abseta$ bin.

\begin{figure}[htbp]
\centering
\includegraphics[width=0.92\textwidth]{llgamma/baseline/plots/rew_reta.pdf}
\caption{Cell-computed \Reta\ comparison: data versus uncorrected MC,
  M1 (shift only), and M2 (shift+stretch), for the combined channel,
  unconverted photons.}
\label{fig:rew-reta}
\end{figure}

\begin{figure}[htbp]
\centering
\includegraphics[width=0.92\textwidth]{llgamma/baseline/plots/rew_rphi.pdf}
\caption{Cell-computed \Rphi\ comparison (same layout as
  Figure~\ref{fig:rew-reta}).}
\label{fig:rew-rphi}
\end{figure}

\begin{figure}[htbp]
\centering
\includegraphics[width=0.92\textwidth]{llgamma/baseline/plots/rew_weta2.pdf}
\caption{Cell-computed \weta\ comparison (same layout as
  Figure~\ref{fig:rew-reta}).}
\label{fig:rew-weta2}
\end{figure}

\subsection{Chi-squared comparison}
\label{sec:results-chi2}

Table~\ref{tab:chi2-llgamma-baseline} reports the $\chisq$ between each
MC variant and data, computed bin-by-bin after area normalisation, for
each $\abseta$ bin.  The Pearson $\chi^2$ is defined as
\begin{equation}
  \chi^2 = \sum_{b} \frac{(h^{\text{MC}}_b - h^{\text{data}}_b)^2}
                         {h^{\text{data}}_b} \,,
\end{equation}
where the sum runs over histogram bins with $h^{\text{data}}_b > 0$.

\subsection{Average performance}
\label{sec:results-average}

Table~\ref{tab:chi2-summary} summarises the average $\chisq$ across all
non-empty $\abseta$ bins for each method, channel, and conversion
scenario.

% Auto-generated by extract_chi2.C — do not edit by hand
\begin{table}[htbp]
\centering
\caption{Average $\chisq$ across all non-empty $|\eta|$ bins for each correction method.  Lower is better.}
\label{tab:chi2-summary}
\resizebox{\textwidth}{!}{%
\begin{tabular}{ll|rrrr|rrrr|rrrr}
\toprule
 & & \multicolumn{4}{c|}{$R_{\eta}$} & \multicolumn{4}{c|}{$R_{\phi}$} & \multicolumn{4}{c}{$w_{\eta 2}$} \\
Channel & Scenario & MC & M1 & M2 & Fudged & MC & M1 & M2 & Fudged & MC & M1 & M2 & Fudged \\
\midrule
$Z \to \ell\ell\gamma$ & Unconverted & 0.0015 & 0.0006 & 0.0005 & 0.0005 & 0.0007 & 0.0011 & 0.0006 & 0.0005 & 0.0076 & 0.0007 & 0.0005 & 0.0005 \\
$Z \to \ell\ell\gamma$ & Converted & 0.0020 & 0.0013 & 0.0015 & 0.0010 & 0.0010 & 0.0010 & 0.0010 & 0.0010 & 0.0051 & 0.0011 & 0.0010 & 0.0009 \\
$Z \to \ell\ell\gamma$ & Inclusive & 0.0012 & 0.0004 & 0.0004 & 0.0003 & 0.0005 & 0.0006 & 0.0004 & 0.0003 & 0.0055 & 0.0005 & 0.0004 & 0.0003 \\
\bottomrule
\end{tabular}}
\end{table}



%% ====================================================================
\section{Conversion Scenarios}
\label{sec:conversion}
%% ====================================================================

The corrections are re-derived independently for each conversion
scenario.  Converted photons have different cluster energy profiles
due to the upstream conversion, and the correction parameters adapt
accordingly.

\subsection{Converted photons}
\label{sec:conv-converted}

Figure~\ref{fig:rew-reta-converted} shows the \Reta\ comparison for
converted photons in the combined channel.  The corresponding $\chisq$
values are reported in Table~\ref{tab:chi2-llgamma-converted}
(Appendix~\ref{app:chi2-tables}).

\begin{figure}[htbp]
\centering
\includegraphics[width=0.92\textwidth]{llgamma/converted/plots/rew_reta.pdf}
\caption{Cell-computed \Reta\ for the combined channel, converted
  photons: data versus MC, M1, and M2.}
\label{fig:rew-reta-converted}
\end{figure}

\subsection{All conversions combined}
\label{sec:conv-allconv}

Figure~\ref{fig:rew-reta-allconv} shows the \Reta\ comparison without
any conversion requirement.  Full $\chisq$ tables are in
Table~\ref{tab:chi2-llgamma-all_conv} (Appendix~\ref{app:chi2-tables}).

\begin{figure}[htbp]
\centering
\includegraphics[width=0.92\textwidth]{llgamma/all_conv/plots/rew_reta.pdf}
\caption{Cell-computed \Reta\ for the combined channel, all conversions:
  data versus MC, M1, and M2.}
\label{fig:rew-reta-allconv}
\end{figure}

\subsection{Comparison across scenarios}
\label{sec:conv-comparison}

The summary table (Table~\ref{tab:chi2-summary}) allows direct
comparison of the average $\chisq$ across conversion scenarios.

Unconverted photons consistently produce the lowest $\chisq$ values for
all methods.  For the combined channel:\ M1 achieves an average \weta\
$\chisq$ of 0.0007 (unconverted) versus 0.0012 (converted).  Converted
photons produce more diffuse electromagnetic showers due to the upstream
$e^+e^-$ pair production, leading to broader cell-fraction distributions
whose correction is intrinsically harder.  The reduced data statistics
for converted photons (45\,532 events versus 205\,309 unconverted in the
combined channel) further limit the precision of the derived corrections.

The inclusive (all-conversions) scenario contains both populations and
benefits from the highest statistics; it generally falls between the
unconverted and converted performance.


%% ====================================================================
\section{Per-Channel Results}
\label{sec:per-channel}
%% ====================================================================

The corrections are derived and evaluated separately for each decay
channel.  Tables for the $Z \to ee\gamma$ and $Z \to \mu\mu\gamma$
channels are available in Appendix~\ref{app:chi2-tables}
(Tables~\ref{tab:chi2-eegamma-baseline}--\ref{tab:chi2-mumugamma-all_conv}).

Both channels show consistent behaviour: the correction performance is
stable across the two lepton flavours, validating the use of the
combined channel as the primary result.  The $Z \to \mu\mu\gamma$
channel yields $\sim$60\% more data events (127\,178 versus 78\,131
unconverted) and correspondingly lower $\chisq$ values in all
scenarios, confirming that increased statistics improve correction
quality.  No channel-specific pathologies are observed; the M2
instabilities discussed in Section~\ref{sec:disc-m2-instability} appear
in different $\abseta$ bins in the individual channels but are present
in both, confirming that they are a property of the method rather than
channel-specific artefacts.


%% ====================================================================
\section{Cell Energy Profile Maps}
\label{sec:cell-profiles}
%% ====================================================================

The pipeline produces $7 \times 11$ heatmaps of the mean cell energy
fraction for data, uncorrected MC, and corrected MC, as well as
residual (delta) maps and correction vectors.

\subsection{Before and after correction}

Figure~\ref{fig:cell-before} shows the mean cell fraction profiles for
data and uncorrected MC.  The difference $\Delta_k = \langle
f_{\text{data}} \rangle_k - \langle f_{\text{MC}} \rangle_k$ is the
M1 correction.

\begin{figure}[htbp]
\centering
\includegraphics[width=0.48\textwidth]{llgamma/baseline/plots/cell_data.pdf}%
\hfill
\includegraphics[width=0.48\textwidth]{llgamma/baseline/plots/cell_mc.pdf}
\caption{Mean cell energy fraction profiles before correction:
  data (left) and MC (right), for each $\abseta$ bin.}
\label{fig:cell-before}
\end{figure}

After M1 correction the MC profile matches data by construction.
M2 additionally adjusts the cell fraction variance.  Both corrected
profiles are shown in Figure~\ref{fig:cell-after-m1}.

\begin{figure}[htbp]
\centering
\includegraphics[width=0.48\textwidth]{llgamma/baseline/plots/cell_mc_m1.pdf}%
\hfill
\includegraphics[width=0.48\textwidth]{llgamma/baseline/plots/cell_mc_m2.pdf}
\caption{Mean cell fraction after correction: M1 (left) and M2 (right),
  for each $\abseta$ bin.  Compare with data in
  Figure~\ref{fig:cell-before}.}
\label{fig:cell-after-m1}
\end{figure}

\subsection{M2 correction vectors}

The M2 shift and stretch parameters are visualised as 1D vectors per
$\abseta$ bin in Figures~\ref{fig:cell-shift}~and~\ref{fig:cell-stretch}.

\begin{figure}[htbp]
\centering
\includegraphics[width=0.92\textwidth]{llgamma/baseline/plots/cell_shift.pdf}
\caption{M2 shift correction vector $\text{shift}_k$ for each cell~$k$,
  per $\abseta$ bin.}
\label{fig:cell-shift}
\end{figure}

\begin{figure}[htbp]
\centering
\includegraphics[width=0.92\textwidth]{llgamma/baseline/plots/cell_stretch.pdf}
\caption{M2 stretch correction vector $\text{stretch}_k$ for each
  cell~$k$, per $\abseta$ bin.}
\label{fig:cell-stretch}
\end{figure}


%% ====================================================================
\section{Discussion}
\label{sec:discussion}
%% ====================================================================

The nine-scenario evaluation reveals a clear hierarchy among the
correction approaches.  This section presents a quantitative analysis,
diagnoses the M2 instabilities, and compares the cell-level methods
with the standard ATLAS fudge-factor corrections.

\subsection{M1 performance across all scenarios}
\label{sec:disc-m1}

Method~1 delivers consistent and substantial improvements across all
nine channel$\times$scenario combinations.  For the primary benchmark
($Z \to \ell\ell\gamma$, unconverted):
\begin{itemize}
  \item \Reta: average $\chisq$ reduced from $0.0015$ (uncorrected MC)
        to $0.0006$ (M1), a factor~2.5 improvement.
  \item \weta: reduced from $0.0076$ to $0.0007$, a factor~10.9
        improvement --- the largest gain of any method for any variable.
  \item \Rphi: increases slightly from $0.0007$ to $0.0011$.
\end{itemize}
The pattern is reproduced across all channels and conversion scenarios
(Table~\ref{tab:chi2-summary}).  M1's strength for \Reta\ and \weta\
is expected: both variables are dominated by the central cells of the
$7 \times 11$ cluster, where the mean cell fraction is large and the
shift correction has the most leverage.

The \Rphi\ anomaly merits discussion.  $R_\phi = E(3\!\times\!3) /
E(3\!\times\!7)$ is a ratio of two overlapping energy sums sharing the
same $3\!\times\!3$ core.  The M1 shift corrects each cell
independently; while the individual cells improve in mean, the
\emph{ratio} of two correlated sums need not improve --- the correction
introduces coherent shifts in both numerator and denominator that can
increase the spread of the ratio distribution.  This is a fundamental
limitation of cell-level mean corrections for ratio observables whose
numerator and denominator share cells.  In practice the degradation is
modest ($0.0007 \to 0.0011$) and M1's \Rphi\ $\chisq$ remains well
below the uncorrected \weta\ value, so the overall correction quality
is still acceptable.

M1 preserves total cluster energy exactly by construction
($\sum_k \Delta_k = 0$), requiring no rescaling.  This property makes
M1 robust against pathological cell configurations and suitable for
production use without additional safeguards.

\subsection{M2 instability diagnosis}
\label{sec:disc-m2-instability}

Contrary to the expectation that matching both the mean and variance
of the cell fraction should uniformly improve agreement, M2 exhibits
\textbf{catastrophic failures} in specific $\abseta$ bins.
Table~\ref{tab:m2-failures} lists the most severe cases.

\begin{table}[htbp]
\centering
\caption{M2 catastrophic failures: $\abseta$ bins where M2 produces
  $\chisq$ values more than $10\times$ worse than uncorrected MC.
  Only the worst instances per scenario are shown; additional bins
  with $5$--$10\times$ degradation exist in most scenarios.}
\label{tab:m2-failures}
\begin{tabular}{llllrrl}
\toprule
Channel & Scenario & $\abseta$ bin & Variable
  & $\chisq^{\text{MC}}$ & $\chisq^{\text{M2}}$
  & Ratio \\
\midrule
$ee\gamma$ & Unconv.
  & {[}0.60, 0.80) & $R_\phi$ & 0.0004 & 0.0323 & $81\times$ \\
  & & {[}0.60, 0.80) & $w_{\eta 2}$ & 0.0017 & 0.0481 & $28\times$ \\
  & & {[}1.20, 1.30) & $R_\phi$ & 0.0011 & 0.0641 & $58\times$ \\
  & & {[}1.20, 1.30) & $w_{\eta 2}$ & 0.0016 & 0.0273 & $17\times$ \\
\midrule
$\mu\mu\gamma$ & Converted
  & {[}0.20, 0.40) & $R_\phi$ & 0.0012 & 0.0426 & $36\times$ \\
  & & {[}0.20, 0.40) & $w_{\eta 2}$ & 0.0033 & 0.0463 & $14\times$ \\
  & & {[}1.00, 1.20) & $R_\phi$ & 0.0004 & 0.0360 & $90\times$ \\
  & & {[}1.00, 1.20) & $w_{\eta 2}$ & 0.0011 & 0.0588 & $53\times$ \\
\midrule
$\ell\ell\gamma$ & Converted
  & {[}0.20, 0.40) & $R_\phi$ & 0.0007 & 0.0262 & $37\times$ \\
  & & {[}0.20, 0.40) & $w_{\eta 2}$ & 0.0025 & 0.0509 & $20\times$ \\
  & & {[}1.00, 1.20) & $R_\phi$ & 0.0003 & 0.0238 & $79\times$ \\
  & & {[}1.00, 1.20) & $w_{\eta 2}$ & 0.0008 & 0.0159 & $20\times$ \\
\bottomrule
\end{tabular}
\end{table}

\paragraph{Root cause: systematically large stretch factors in peripheral cells.}
The M2 stretch factor is
$\text{stretch}_k = \sigma_{\text{data},k} / \sigma_{\text{MC},k}$.
Inspection of the computed stretch vectors reveals a striking
asymmetry: stretch values are almost exclusively $\geq 1$, and often
\textbf{massively} so.  The minimum stretch across all 77~cells and
14~$\abseta$ bins is $\sim$0.93 --- barely below unity --- while the
maximum reaches $\sim$16 (\texttt{eegamma} baseline, $\abseta$ bin~3,
cell~44) and $\sim$10 (\texttt{llgamma} baseline, same bin/cell).

This means $\sigma_{\text{data},k} > \sigma_{\text{MC},k}$
\textbf{essentially everywhere}: data cell-energy distributions are
systematically wider than MC.  The likely cause is that MC showers are
``too clean'' --- they lack the full spectrum of detector noise,
pile-up fluctuations, and material-interaction tails that broaden the
real data distributions.

The magnitude of the discrepancy depends on the energy fraction
$\langle f_k \rangle$ carried by the cell.  For the central cells
($\langle f_k \rangle \sim 0.02$--$0.10$), both $\sigma_{\text{data}}$
and $\sigma_{\text{MC}}$ are sizable and their ratio stays in the range
0.95--1.3 --- benign corrections.  For peripheral cells
($\langle f_k \rangle < 10^{-3}$), both standard deviations are tiny
($\sim$$10^{-4}$), and even a small absolute difference
$\sigma_{\text{data}} - \sigma_{\text{MC}}$ translates into a ratio
$\gg 1$.  In the catastrophic $\abseta$ bins:
\begin{itemize}
  \item $\abseta \in [0.60, 0.80)$ (bin~3, unconverted): 18 of 77 cells
    have $\text{stretch} > 2$, with a maximum of 10.2;
  \item $\abseta \in [1.20, 1.30)$ (bin~6, unconverted): 30 of 77 cells
    have $\text{stretch} > 2$, with a maximum of 9.8;
  \item for the \texttt{eegamma} channel (lower statistics), the maxima
    rise to 16.3 and 15.3 in the same bins.
\end{itemize}

These large stretch factors \textbf{amplify} the cell-energy
fluctuations far beyond their physical range.  After the
energy-conservation rescaling~\eqref{eq:m2-rescale}, the amplified
peripheral cells steal energy from the central core, distorting the
shower shape variables and producing the observed catastrophic $\chisq$
values.

Crucially, the reverse problem ($\text{stretch} \ll 1$) is
\textbf{not observed}: MC never significantly over-estimates the
cell-energy variance.  The instability is therefore one-sided: it is
always an amplification, never a compression.

This is \textbf{not a bug} but an \textbf{intrinsic limitation} of the
variance-matching approach when applied to high-dimensional clusters
where most cells carry negligible energy and the MC systematically
underestimates the cell-level variance.  The problem is exacerbated
for converted photons, which have more diffuse showers and hence more
cells with intermediate energy fractions in the regime where the
variance ratio is most unstable.

Notably, the failing $\abseta$ bins differ between unconverted
({[}0.60, 0.80) and {[}1.20, 1.30)) and converted ({[}0.20, 0.40) and
{[}1.00, 1.20)) photons.  This bin-dependent pattern further supports
the diagnosis: the instability depends on the specific shower topology
in each ($\abseta$, conversion type) combination, rather than on a fixed
detector geometry effect.

\paragraph{M2 averaged performance.}
Table~\ref{tab:chi2-summary} shows that M2 \textbf{never outperforms
M1 on average} across any of the 27 (channel $\times$ scenario $\times$
variable) combinations.  The catastrophic bins dominate the average,
overwhelming the improvements that M2 provides in other bins.  For
example, in the combined channel (unconverted), M2 achieves excellent
$\chisq$ in the high-$\abseta$ region ($\abseta > 1.8$), often matching
or surpassing M1 --- but this is erased by a single bin where $\chisq$
is $50\times$ worse.

\subsection{Comparison with fudge factors}
\label{sec:disc-fudge}

The ATLAS fudge-factor framework applies an empirical, per-variable
correction to the branch-level shower shape values.  The ``Fudged''
column in the $\chisq$ tables represents the agreement between data
and MC after these corrections.

Fudge factors achieve the lowest average $\chisq$ in the majority of
the 27 (channel $\times$ scenario $\times$ variable) combinations,
typically outperforming M1 by 20--40\%.  For instance, in
$Z \to \ell\ell\gamma$ (unconverted):\ \Reta\ $0.0005$ (Fudged) vs.\
$0.0006$ (M1); \weta\ $0.0005$ vs.\ $0.0007$.

This marginal advantage is expected: fudge factors are tuned
\emph{independently per variable} using the full data--MC comparison for
that specific observable, whereas cell-level methods derive a
\emph{single} set of 77--cell corrections from which all shower shapes
are recomputed.  The cell-level approach has fewer effective degrees of
freedom (77~shifts per $\abseta$ bin) to simultaneously correct three
observables built from overlapping cell windows.

The key advantage of cell-level corrections is \textbf{physics
generality}: any observable built from the Layer\,2 cell energies ---
including variables not currently used in photon ID such as
$E_{\text{ratio}}$, $f_{\text{side}}$, or new discrimination variables
--- benefits automatically from the same corrections without separate
tuning.  Fudge factors would need to be re-derived for each new variable.

A second advantage is \textbf{correlation preservation}: because the
cell energies are corrected before the shower shapes are computed, the
correlations between \Reta, \Rphi, and \weta\ are preserved by
construction.  Fudge factors, applied independently per variable, can
distort these correlations.

In summary, M1 is the recommended cell-level method.  It provides
\emph{most} of the fudge-factor improvement (within $\sim$30\%) while
offering a more physically motivated and extensible correction
framework.  Fudge factors remain the gold standard for per-variable
agreement when only the three standard shower shapes are of interest.

\subsection{Problematic \texorpdfstring{$\abseta$}{|eta|} regions}
\label{sec:disc-problematic}

Several $\abseta$ regions exhibit elevated $\chisq$ for all methods:

\paragraph{High $\abseta$ ($> 2.0$).}
The uncorrected MC shows the largest mismodelling for \weta\ in the
forward region: $\chisq = 0.027$ at $\abseta \in [2.0, 2.2)$ and
$\chisq = 0.030$ at $\abseta \in [2.2, 2.4)$ for the combined
unconverted channel.  M1 reduces these by factors of 50--60 to
$\chisq \approx 0.0005$, demonstrating that the cell-fraction mean
shift captures the dominant mismodelling even where the detector
geometry differs significantly from the barrel.

\paragraph{Near-crack region ($\abseta \in [1.30, 1.52)$).}
The bins immediately adjacent to the calorimeter crack show elevated
$\chisq$ for all methods.  The $[1.30, 1.37)$ bin, in particular,
has limited data statistics (704 events for $ee\gamma$ unconverted)
and residual material effects from the crack region that affect the
shower profile.  All methods produce $\chisq \sim 0.01$--$0.03$ here,
which is 5--10$\times$ worse than the barrel average.  The
$[1.52, 1.60)$ bin, immediately after the crack, shows a similar
pattern.  These bins are unlikely to improve without dedicated
treatment (e.g.\ separate correction parameters for the transition
region).

\paragraph{Converted photons.}
All methods degrade for converted photons relative to unconverted, with
the degradation factor roughly proportional to $\sqrt{N_{\text{unconv}} /
N_{\text{conv}}}$, suggesting that statistical precision of the
corrections is a limiting factor.  For the combined channel, the
converted dataset has 45\,532 events versus 205\,309 unconverted (a
factor~4.5 fewer), corresponding to a $\sim$$2\times$ increase in
statistical uncertainty on the correction parameters.


%% ====================================================================
\section{Future Work}
\label{sec:future}
%% ====================================================================

\subsection{M2 regularisation}
\label{sec:future-regularisation}

The M2 instabilities identified in Section~\ref{sec:disc-m2-instability}
can be mitigated by regularising the stretch factor.  The key empirical
observation is that the problem is \textbf{one-sided}: stretch values
are almost exclusively $\geq 1$ (data distributions are systematically
wider than MC), with extremes reaching $\sim$10--16 in the worst bins,
while the minimum observed stretch is $\sim$0.93.  Regularisation
therefore needs primarily to cap the upper tail.  Several strategies
are possible:

\begin{enumerate}
  \item \textbf{Winsorisation (upper cap)}: restrict the stretch factor to
        $\text{stretch}_k \leq s_{\max}$, e.g.\ $s_{\max} = 1.5$ or
        $2.0$.  Since the instability is exclusively upward, a lower
        bound is not needed in practice ($s_{\min} = 1.0$ suffices).
        This is the simplest approach and guarantees bounded corrections,
        at the cost of incomplete variance matching in capped cells.
  \item \textbf{Peripheral-cell exclusion}: apply M2 only to cells with
        $\langle f_k \rangle > f_{\min}$ (e.g.\ $f_{\min} = 0.002$)
        and fall back to M1 for the remaining cells.  This preserves
        M2's variance matching where it is stable while avoiding the
        problematic near-zero regime.
  \item \textbf{Regularised variance}: replace $\sigma_{\text{MC},k}$
        with
        $\sigma_{k,\text{reg}} = \sqrt{\sigma_{\text{MC},k}^2
        + \varepsilon^2}$ for a small regularisation constant
        $\varepsilon$, preventing the ratio from diverging when the
        MC variance is small.  The optimal $\varepsilon$ can be
        chosen by minimising the $\chisq$ on a held-out validation
        sample.
\end{enumerate}

The peripheral-cell exclusion approach is the most physically motivated:
cells with $f_k < 10^{-3}$ contribute negligibly to all standard shower
shapes, and their variance carries no useful physics information.  A
systematic study of $f_{\min}$ thresholds, evaluating $\chisq$ across
all $\abseta$ bins, is the recommended first step.

\subsection{\texorpdfstring{$\pt$}{pT}-dependent corrections}
\label{sec:future-pt}

The current corrections are derived in 14~$\abseta$ bins with no
dependence on photon \pt.  Francisco Almeida's \texttt{photoncellbasedrw}
framework uses $6~\pt \times 5~\text{hit-position} = 30$ bins per
conversion type, with \textbf{no} $\abseta$ binning \cite{francisco}.
These two binning strategies are completely orthogonal, addressing
different physics:

\begin{itemize}
  \item $\abseta$ binning captures detector geometry effects (material
        budget, cell granularity, crack proximity).
  \item \pt\ binning captures energy-scale-dependent shower evolution
        (depth of shower maximum, lateral spread).
  \item Hit-position binning captures sub-cell impact-point effects on
        the cell response.
\end{itemize}

The natural extension is a combined $\pt \times \abseta$ grid.
With $6~\pt \times 14~\abseta = 84$ sets of corrections, the statistical
dilution per bin must be monitored: the current dataset provides
$\sim$$16\,\text{k}$ data events per $\abseta$ bin (unconverted,
combined channel); splitting further into 6~\pt\ bins yields
$\sim$$2.7\,\text{k}$ per cell, which may be marginal for stable
variance estimation (relevant to M2, less so for M1).

A staged approach is recommended:
\begin{enumerate}
  \item First evaluate M1 with $\pt \times \abseta$ to quantify the gain
        from \pt-differential corrections.
  \item If the gain is significant, add hit-position binning and assess
        whether the statistical precision remains adequate.
\end{enumerate}

\subsection{Machine-learning approaches}
\label{sec:future-ml}

The M1 and M2 methods apply independent linear corrections to each of
the 77~cells.  This has two fundamental limitations:
(a)~inter-cell correlations are not preserved beyond the mean (M1) or
first two moments (M2), and (b)~higher-order (non-Gaussian) features of
the cell energy distribution are not captured.  Machine-learning methods
can address both.

\subsubsection{Normalising flows}

A normalising flow \cite{papamakarios} learns a bijective transformation
$g: \mathbb{R}^{77} \to \mathbb{R}^{77}$ such that the
distribution of $g(\mathbf{f}_{\text{MC}})$ matches
$\mathbf{f}_{\text{data}}$, where $\mathbf{f}$ is the 77-dimensional
vector of cell energy fractions.

\begin{itemize}
  \item \textbf{Full multivariate correction}: the transformation
        operates on the joint distribution of all 77~cells
        simultaneously, automatically preserving inter-cell correlations
        to all orders.
  \item \textbf{Conditional flows}: the transformation can be
        conditioned on $(\pt, \eta)$, providing differential corrections
        without explicit binning and the associated statistical
        fragmentation.
  \item \textbf{No variance-ratio instability}: normalising flows do not
        compute explicit $\sigma$ ratios; the transformation is learned
        end-to-end, naturally avoiding the M2 failure mode.
  \item \textbf{Non-Gaussian tails}: normalising flows can model
        arbitrary distributions, including the heavy tails and
        near-zero peaks observed in peripheral cell fractions.
\end{itemize}

The training data would consist of cell fraction vectors from data and MC
events passing the same $Z \to \ell\ell\gamma$ selection, with the flow
trained to transport the MC distribution to data.

\subsubsection{Quantile regression and mapping}

An alternative per-cell approach is quantile mapping: for each cell~$k$,
the MC fraction $f_k^{\text{MC}}$ is mapped to the data distribution via
\begin{equation}
  f_k^{\text{corr}} = F_{\text{data},k}^{-1}\bigl(
    F_{\text{MC},k}(f_k^{\text{MC}})\bigr) \,,
\end{equation}
where $F$ denotes the cumulative distribution function.  This can be
implemented empirically (histogram-based CDF) or learned via
neural-network quantile regression \cite{koenker}.

Quantile mapping is non-linear and can correct arbitrary distributional
differences on a per-cell basis, improving on the linear M1/M2
transformations.  It can also be extended to conditional quantile
mapping, where $F$ depends on $(\pt, \eta)$.  However, like M1/M2, it
operates per cell and does not intrinsically preserve inter-cell
correlations.

\subsubsection{Practical considerations}

Machine-learning approaches offer the theoretical promise of superior
corrections but face practical challenges:
\begin{itemize}
  \item \textbf{Training data}: the $Z \to \ell\ell\gamma$ sample
        provides $\sim$$200\,\text{k}$ data events and
        $\sim$$2.8\,\text{M}$ MC events (unconverted, combined channel).
        This may be limiting for 77-dimensional flows.
  \item \textbf{Overfitting}: regularisation, early stopping, and
        validation splits are essential; the corrections must be
        evaluated on a statistically independent dataset.
  \item \textbf{Inference in C++}: ATLAS production workflows require
        fast C++ inference, necessitating ONNX export or equivalent.
  \item \textbf{Transferability}: corrections derived on FSR photons
        must be validated for other photon populations (e.g.\ $H \to
        \gamma\gamma$, $\gamma$+jets), where the \pt\ spectrum and
        isolation environment differ.
\end{itemize}

Given the M2 instability results, normalising flows represent the most
promising direction: they avoid explicit moment matching, preserve the
full multivariate structure, and can be conditioned on kinematics
continuously.


%% ====================================================================
\section{Summary}
\label{sec:summary}
%% ====================================================================

This note documents a cell-level energy reweighting pipeline for
correcting photon shower shape variables (\Reta, \Rphi, \weta) in ATLAS
Run\,3 simulation.  Using $Z \to \ell\ell\gamma$ events from 2024 data
and MC23e, corrections are derived and applied in 14~$\abseta$ bins
across three decay channels and three conversion scenarios (nine
configurations total, $\sim$$250\,\text{k}$ data events in the
combined unconverted sample).

Two methods are implemented:
\begin{itemize}
  \item \textbf{M1 (flat shift)}: corrects the mean cell energy fraction
        with exact energy conservation.  Reduces the average $\chisq$
        by factors of 2--11 across all scenarios.  M1 is the
        \textbf{recommended cell-level method}: it is stable in every
        $\abseta$ bin, requires no regularisation, and achieves
        performance within $\sim$30\% of the empirical fudge factors.
  \item \textbf{M2 (shift+stretch)}: additionally matches the cell
        fraction variance.  \textbf{Not recommended} in its current form
        due to catastrophic instabilities in specific $\abseta$ bins
        (up to $90\times$ worse than uncorrected MC), caused by
        divergent variance ratios in peripheral cells.  Regularisation
        of the stretch factor (Section~\ref{sec:future-regularisation})
        is required before M2 can be used reliably.
\end{itemize}

Fudge factors remain the most precise per-variable corrections,
marginally outperforming M1.  However, cell-level corrections offer
two distinct advantages: (a)~a single set of corrections simultaneously
improves \emph{all} observables built from cell energies, and (b)~the
inter-variable correlations are preserved by construction.

Priority directions for future work include regularisation of M2
(Section~\ref{sec:future-regularisation}), $\pt$-differential
corrections (Section~\ref{sec:future-pt}), and normalising flows
that avoid explicit moment matching entirely
(Section~\ref{sec:future-ml}).


%% ====================================================================
\appendix
\section{Per-Eta Chi-Squared Tables}
\label{app:chi2-tables}
%% ====================================================================

This appendix contains the per-$\abseta$-bin $\chisq$ tables for all
channel and conversion scenario combinations.

% Auto-generated by extract_chi2.C — filtered to llgamma only
% Per-$|\eta|$-bin $\chi^2/n_{\mathrm{df}}$ tables for $Z \to \ell\ell\gamma$.

\begin{table}[htbp]
\centering
\caption{Per-$|\eta|$-bin $\chisq$ for $Z \to \ell\ell\gamma$, Unconverted scenario.  Columns show $\chisq$ between each MC variant and data (cell-computed).}
\label{tab:chi2-llgamma-unconverted}
\resizebox{\textwidth}{!}{%
\begin{tabular}{r|rrrr|rrrr|rrrr}
\toprule
$|\eta|$ bin & \multicolumn{4}{c|}{$R_{\eta}$} & \multicolumn{4}{c|}{$R_{\phi}$} & \multicolumn{4}{c}{$w_{\eta 2}$} \\
(range) & MC & M1 & M2 & Fud & MC & M1 & M2 & Fud & MC & M1 & M2 & Fud \\
\midrule
{[}0.00, 0.20) & 0.0002 & 0.0002 & 0.0002 & 0.0002 & 0.0004 & 0.0005 & 0.0003 & 0.0002 & 0.0008 & 0.0003 & 0.0002 & 0.0003 \\
{[}0.20, 0.40) & 0.0002 & 0.0001 & 0.0001 & 0.0001 & 0.0003 & 0.0004 & 0.0003 & 0.0002 & 0.0014 & 0.0004 & 0.0003 & 0.0003 \\
{[}0.40, 0.60) & 0.0003 & 0.0002 & 0.0002 & 0.0002 & 0.0005 & 0.0005 & 0.0003 & 0.0002 & 0.0016 & 0.0003 & 0.0003 & 0.0003 \\
{[}0.60, 0.80) & 0.0003 & 0.0002 & 0.0002 & 0.0002 & 0.0003 & 0.0004 & 0.0003 & 0.0002 & 0.0016 & 0.0004 & 0.0003 & 0.0003 \\
{[}0.80, 1.00) & 0.0005 & 0.0002 & 0.0002 & 0.0002 & 0.0007 & 0.0005 & 0.0003 & 0.0003 & 0.0006 & 0.0004 & 0.0002 & 0.0003 \\
{[}1.00, 1.20) & 0.0004 & 0.0001 & 0.0002 & 0.0002 & 0.0010 & 0.0007 & 0.0004 & 0.0003 & 0.0005 & 0.0006 & 0.0003 & 0.0003 \\
{[}1.20, 1.30) & 0.0007 & 0.0004 & 0.0005 & 0.0004 & 0.0008 & 0.0010 & 0.0004 & 0.0004 & 0.0010 & 0.0007 & 0.0004 & 0.0004 \\
{[}1.30, 1.37) & 0.0027 & 0.0011 & 0.0017 & 0.0019 & 0.0014 & 0.0016 & 0.0011 & 0.0014 & 0.0022 & 0.0020 & 0.0019 & 0.0012 \\
{[}1.52, 1.60) & 0.0014 & 0.0027 & 0.0015 & 0.0008 & 0.0010 & 0.0036 & 0.0011 & 0.0011 & 0.0072 & 0.0014 & 0.0011 & 0.0010 \\
{[}1.60, 1.80) & 0.0013 & 0.0005 & 0.0004 & 0.0003 & 0.0004 & 0.0011 & 0.0006 & 0.0005 & 0.0080 & 0.0004 & 0.0004 & 0.0003 \\
{[}1.80, 2.00) & 0.0023 & 0.0005 & 0.0009 & 0.0004 & 0.0005 & 0.0011 & 0.0009 & 0.0004 & 0.0171 & 0.0015 & 0.0008 & 0.0004 \\
{[}2.00, 2.20) & 0.0034 & 0.0004 & 0.0004 & 0.0006 & 0.0006 & 0.0009 & 0.0005 & 0.0004 & 0.0273 & 0.0005 & 0.0003 & 0.0004 \\
{[}2.20, 2.40) & 0.0054 & 0.0006 & 0.0007 & 0.0005 & 0.0011 & 0.0014 & 0.0007 & 0.0008 & 0.0296 & 0.0005 & 0.0005 & 0.0005 \\
\midrule
Average & 0.0015 & 0.0006 & 0.0005 & 0.0005 & 0.0007 & 0.0011 & 0.0006 & 0.0005 & 0.0076 & 0.0007 & 0.0005 & 0.0005 \\
\bottomrule
\end{tabular}}
\end{table}

\begin{table}[htbp]
\centering
\caption{Per-$|\eta|$-bin $\chisq$ for $Z \to \ell\ell\gamma$, Converted scenario.  Columns show $\chisq$ between each MC variant and data (cell-computed).}
\label{tab:chi2-llgamma-converted}
\resizebox{\textwidth}{!}{%
\begin{tabular}{r|rrrr|rrrr|rrrr}
\toprule
$|\eta|$ bin & \multicolumn{4}{c|}{$R_{\eta}$} & \multicolumn{4}{c|}{$R_{\phi}$} & \multicolumn{4}{c}{$w_{\eta 2}$} \\
(range) & MC & M1 & M2 & Fud & MC & M1 & M2 & Fud & MC & M1 & M2 & Fud \\
\midrule
{[}0.00, 0.20) & 0.0021 & 0.0023 & 0.0032 & 0.0022 & 0.0014 & 0.0019 & 0.0017 & 0.0015 & 0.0021 & 0.0016 & 0.0022 & 0.0018 \\
{[}0.20, 0.40) & 0.0009 & 0.0013 & 0.0017 & 0.0010 & 0.0007 & 0.0012 & 0.0024 & 0.0011 & 0.0025 & 0.0041 & 0.0017 & 0.0017 \\
{[}0.40, 0.60) & 0.0011 & 0.0015 & 0.0022 & 0.0016 & 0.0013 & 0.0014 & 0.0014 & 0.0011 & 0.0020 & 0.0010 & 0.0010 & 0.0011 \\
{[}0.60, 0.80) & 0.0013 & 0.0033 & 0.0029 & 0.0010 & 0.0009 & 0.0009 & 0.0010 & 0.0009 & 0.0017 & 0.0010 & 0.0011 & 0.0010 \\
{[}0.80, 1.00) & 0.0009 & 0.0007 & 0.0008 & 0.0006 & 0.0006 & 0.0005 & 0.0005 & 0.0005 & 0.0007 & 0.0006 & 0.0005 & 0.0005 \\
{[}1.00, 1.20) & 0.0006 & 0.0005 & 0.0005 & 0.0004 & 0.0003 & 0.0004 & 0.0004 & 0.0006 & 0.0008 & 0.0005 & 0.0007 & 0.0005 \\
{[}1.20, 1.30) & 0.0009 & 0.0007 & 0.0007 & 0.0006 & 0.0008 & 0.0006 & 0.0006 & 0.0008 & 0.0011 & 0.0005 & 0.0006 & 0.0006 \\
{[}1.30, 1.37) & 0.0029 & 0.0023 & 0.0023 & 0.0014 & 0.0023 & 0.0019 & 0.0023 & 0.0016 & 0.0015 & 0.0016 & 0.0017 & 0.0015 \\
{[}1.52, 1.60) & 0.0017 & 0.0018 & 0.0017 & 0.0011 & 0.0007 & 0.0009 & 0.0008 & 0.0008 & 0.0056 & 0.0010 & 0.0006 & 0.0009 \\
{[}1.60, 1.80) & 0.0011 & 0.0003 & 0.0003 & 0.0003 & 0.0004 & 0.0006 & 0.0003 & 0.0006 & 0.0052 & 0.0005 & 0.0004 & 0.0004 \\
{[}1.80, 2.00) & 0.0028 & 0.0007 & 0.0013 & 0.0005 & 0.0005 & 0.0005 & 0.0005 & 0.0007 & 0.0091 & 0.0004 & 0.0004 & 0.0005 \\
{[}2.00, 2.20) & 0.0042 & 0.0008 & 0.0009 & 0.0006 & 0.0008 & 0.0007 & 0.0006 & 0.0006 & 0.0177 & 0.0010 & 0.0008 & 0.0005 \\
{[}2.20, 2.40) & 0.0051 & 0.0013 & 0.0012 & 0.0012 & 0.0017 & 0.0013 & 0.0008 & 0.0017 & 0.0164 & 0.0012 & 0.0011 & 0.0009 \\
\midrule
Average & 0.0020 & 0.0013 & 0.0015 & 0.0010 & 0.0010 & 0.0010 & 0.0010 & 0.0010 & 0.0051 & 0.0011 & 0.0010 & 0.0009 \\
\bottomrule
\end{tabular}}
\end{table}

\begin{table}[htbp]
\centering
\caption{Per-$|\eta|$-bin $\chisq$ for $Z \to \ell\ell\gamma$, Inclusive scenario.  Columns show $\chisq$ between each MC variant and data (cell-computed).}
\label{tab:chi2-llgamma-inclusive}
\resizebox{\textwidth}{!}{%
\begin{tabular}{r|rrrr|rrrr|rrrr}
\toprule
$|\eta|$ bin & \multicolumn{4}{c|}{$R_{\eta}$} & \multicolumn{4}{c|}{$R_{\phi}$} & \multicolumn{4}{c}{$w_{\eta 2}$} \\
(range) & MC & M1 & M2 & Fud & MC & M1 & M2 & Fud & MC & M1 & M2 & Fud \\
\midrule
{[}0.00, 0.20) & 0.0002 & 0.0002 & 0.0002 & 0.0002 & 0.0004 & 0.0005 & 0.0003 & 0.0002 & 0.0008 & 0.0003 & 0.0002 & 0.0003 \\
{[}0.20, 0.40) & 0.0002 & 0.0001 & 0.0001 & 0.0001 & 0.0003 & 0.0003 & 0.0002 & 0.0002 & 0.0013 & 0.0004 & 0.0003 & 0.0003 \\
{[}0.40, 0.60) & 0.0003 & 0.0002 & 0.0002 & 0.0001 & 0.0005 & 0.0004 & 0.0003 & 0.0002 & 0.0015 & 0.0003 & 0.0003 & 0.0003 \\
{[}0.60, 0.80) & 0.0003 & 0.0002 & 0.0002 & 0.0002 & 0.0003 & 0.0004 & 0.0002 & 0.0001 & 0.0014 & 0.0003 & 0.0002 & 0.0003 \\
{[}0.80, 1.00) & 0.0004 & 0.0002 & 0.0002 & 0.0002 & 0.0006 & 0.0004 & 0.0002 & 0.0002 & 0.0005 & 0.0004 & 0.0003 & 0.0003 \\
{[}1.00, 1.20) & 0.0003 & 0.0001 & 0.0002 & 0.0002 & 0.0006 & 0.0005 & 0.0002 & 0.0002 & 0.0004 & 0.0005 & 0.0003 & 0.0003 \\
{[}1.20, 1.30) & 0.0005 & 0.0003 & 0.0003 & 0.0002 & 0.0005 & 0.0006 & 0.0003 & 0.0002 & 0.0008 & 0.0004 & 0.0003 & 0.0004 \\
{[}1.30, 1.37) & 0.0018 & 0.0009 & 0.0012 & 0.0010 & 0.0009 & 0.0011 & 0.0008 & 0.0008 & 0.0014 & 0.0008 & 0.0008 & 0.0007 \\
{[}1.52, 1.60) & 0.0011 & 0.0016 & 0.0011 & 0.0006 & 0.0004 & 0.0013 & 0.0006 & 0.0004 & 0.0053 & 0.0009 & 0.0007 & 0.0006 \\
{[}1.60, 1.80) & 0.0010 & 0.0003 & 0.0002 & 0.0001 & 0.0002 & 0.0003 & 0.0003 & 0.0002 & 0.0054 & 0.0002 & 0.0002 & 0.0002 \\
{[}1.80, 2.00) & 0.0022 & 0.0004 & 0.0007 & 0.0003 & 0.0003 & 0.0005 & 0.0005 & 0.0002 & 0.0115 & 0.0007 & 0.0005 & 0.0002 \\
{[}2.00, 2.20) & 0.0030 & 0.0003 & 0.0004 & 0.0002 & 0.0004 & 0.0006 & 0.0005 & 0.0003 & 0.0202 & 0.0005 & 0.0004 & 0.0003 \\
{[}2.20, 2.40) & 0.0047 & 0.0007 & 0.0007 & 0.0005 & 0.0009 & 0.0008 & 0.0005 & 0.0007 & 0.0211 & 0.0005 & 0.0004 & 0.0004 \\
\midrule
Average & 0.0012 & 0.0004 & 0.0004 & 0.0003 & 0.0005 & 0.0006 & 0.0004 & 0.0003 & 0.0055 & 0.0005 & 0.0004 & 0.0003 \\
\bottomrule
\end{tabular}}
\end{table}



%% ====================================================================
\begin{thebibliography}{9}

\bibitem{francisco}
  F.~Almeida,
  ``Shower shape corrections using cell-level reweighting'',
  ATL-COM-PHYS-2021-640 (2021).

\bibitem{run3note}
  M.~Fernandes da Silva,
  ``Photon shower shape studies with Run\,3 data'',
  ATL-COM-PHYS-2025-662 (2025).

\bibitem{papamakarios}
  G.~Papamakarios, E.~Nalisnick, D.~J.~Rezende, S.~Mohamed, and
  B.~Lakshminarayanan,
  ``Normalizing flows for probabilistic modeling and inference'',
  \emph{J.\ Mach.\ Learn.\ Res.} \textbf{22} (2021) 1--64,
  \href{https://arxiv.org/abs/1912.02762}{arXiv:1912.02762}.

\bibitem{koenker}
  R.~Koenker and G.~Bassett,
  ``Regression quantiles'',
  \emph{Econometrica} \textbf{46} (1978) 33--50.

\end{thebibliography}

\end{document}
